\documentclass[11pt]{article}
\usepackage{amsmath,amssymb,amsthm}
\usepackage{graphicx}
\usepackage{booktabs}
\usepackage{multirow}
\usepackage{float}
\usepackage{caption}
\usepackage{subcaption}
\usepackage{geometry}
\usepackage{enumitem}
\usepackage{hyperref}
\geometry{margin=1in}
\captionsetup{font=small,labelfont=bf}
\newcommand{\elim}{\textsc{Elim}}
\title{Data With the Stars: Inverting Elimination Records to Recover Hidden Fan Votes,\\
Replaying Counterfactual Seasons of \textit{Dancing with the Stars}, and Designing a\\
Fairer Judge--Fan Combination Rule}
\author{Team Reference: 2026 MCM Problem C}
\date{January 2026}

\begin{document}
\maketitle

\begin{abstract}
\noindent
We analyze 34 seasons of \textit{Dancing with the Stars} (421 contestant records,
264 season-week elimination events) to address three linked questions: inferring the
secret weekly fan votes, comparing the show's two historical vote-combination rules,
and evaluating the impact of celebrity characteristics and the judges-save mechanism.
Because fan-vote totals are never released, we formulate the first task as an inverse
problem. Our \textbf{Model 1} (Maximum-Entropy Fan-Vote Inversion) proves feasibility
constructively: \emph{every one of the 264 recorded eliminations is consistent with a
fan-vote vector on the simplex under both the rank and percentage rules} (exact LP
margins: median $\gamma_{\max} = 0.075$, all satisfying the analytic $1/n_s$ bound;
maximum-entropy point estimates with calibrated 90\% Dirichlet-bootstrap intervals whose
half-widths range 0.014--0.06 and vary structurally across weeks). The inferred votes
are informative, not artifacts: calibrated draws reproduce eliminations at a mean rate
of 0.738 versus 0.227 for a structureless null ($\chi^2 = 70{,}243$, $p \approx
10^{-300}$; per-week binomial $p < 0.01$ in 253/264 weeks), and judges alone would have
picked the actual eliminated couple in only 32.2\% of weeks. Our \textbf{Model 2}
(counterfactual season replay) re-runs every season week by week under both combination
rules and the season-28+ judges-save over Monte-Carlo fan-vote draws: the winner flips
under the alternate rule with mean probability 0.489 per season and the top-3 in 63\%,
yet neither rule systematically favors fans (fan-favor index 0.536 RANK vs 0.507 PCT);
rather, the rank rule amplifies large judge--fan disagreements at the tie boundary
--- exactly where the famous controversies live. The replay identifies two distinct
controversy types: Jerry Rice and Bristol Palin are rule-sensitive (their runs largely
vanish under the alternate rule), while Billy Ray Cyrus and Bobby Bones are
rule-robust --- Bones' season-27 victory required only ordinary fan support (inferred
fan share 0.115 vs judge share 0.097) and would survive any mechanism the show has
ever used; the fitted judges-save propensity is $\hat\kappa = 0.17$ [0.05, 0.41].
Our \textbf{Model 3} (hierarchical judge/fan decomposition plus fairness--excitement
optimization) finds that celebrity traits (age, industry) are dominated by couple-level
talent/popularity variation, with no trait reaching 95\% credibility in the judge
equation, and that merit protection versus upset drama trades off steeply and
monotonically with the judge weight (fairness $\approx 0.999$ at $\omega = 1$ versus a
95.7\% upset rate at $\omega = 0$), with no interior optimum. We recommend: retain the
50/50 percentage-based combination, keep and formalize the judges-save, publish
fan-vote uncertainty bands rather than totals, and reject the suspense bonus, which our
simulator could not validate. All headline numbers are traced to released artifacts;
where internal review identified unresolved implementation issues (Models 2--3), the
affected results are explicitly labeled provisional rather than validated.
\end{abstract}

\tableofcontents
\newpage

% ============================================================
\section{Introduction}
% ============================================================

\subsection{Problem Background and Motivation}

\textit{Dancing with the Stars} (DWTS), the American adaptation of the British franchise
\textit{Strictly Come Dancing}, has completed 34 seasons of pairing celebrities with
professional dancers, whose weekly performances are scored by a panel of expert judges
and supplemented by fan votes cast by phone and online. Each week the judge scores and
fan votes are combined into a single ordering; the couple with the lowest combined score
is eliminated. The relative weight assigned to fan votes versus judges has changed
repeatedly across the show's history, and each change has been prompted by a public
``controversy'': Jerry Rice reaching the season-2 final despite repeatedly posting the
lowest judges' scores, Bristol Palin's season-11 third-place finish below a dozen better
dances, and Bobby Bones' season-27 victory despite judge scores that were consistently
weak. These controversies are not anecdotes --- they are policy events that changed the
show's rules, from rank-based combination (seasons 1--2), to percentage combination
(seasons 3--27), to the judges-save mechanism introduced around season 28, in which the
bottom two couples are identified by the combined score and the judges themselves choose
who goes home.

At the center of every one of these questions lies an unobservable quantity: \textbf{the
fan votes themselves}. Fan vote totals are never released and are described by the show
as a closely guarded secret. The only public record is the elimination outcome: one
couple leaves each week, and it is presumably (under the rule in force that season) the
couple with the lowest combined score. This is a textbook \emph{inverse problem}: given
the observed judge scores, the observed eliminations, and a known combination rule,
infer the hidden fan votes that are consistent with the entire historical record.

The inverse framing is what makes the problem tractable and, critically, what makes the
certainty of the estimates quantifiable. Instead of guessing a fan-vote predictor and
hoping it matches eliminations, we can compute the \emph{feasible set} of fan-vote
vectors that reproduce every recorded elimination, choose the least-committal point
inside that set, and propagate the residual uncertainty through every downstream
analysis.

\subsection{Data Description}

The provided dataset \texttt{2026\_MCM\_Problem\_C\_Data.csv} contains one row per
celebrity contestant per season, with the following structure:

\begin{itemize}[itemsep=2pt]
  \item \textbf{Contestant information}: season (1--34), celebrity name and partner,
        celebrity demographics including age, gender, and industry/occupation, and the
        assigned professional dancer.
  \item \textbf{Weekly judge scores}: up to 11 weekly columns giving the total judge
        score of that couple in weeks 1--11 of competition (three to eight judges
        depending on season and week, scores on a 0--10 per-judge scale; totals range
        from 8 to 45 in the parsed data).
  \item \textbf{Results}: the elimination/finale outcome for each couple (winner,
        runner-up, week of elimination, or a sentinel for withdrew/forfeit events).
\end{itemize}

Parsing and sanity checks established that the 34 seasons contain \textbf{421
contestant-season records}, that contestant surviving counts decrease monotonically to
each season's finale (finale weeks range from 4 in season 1 to 11 in the longest seasons,
and season 27 opened with 13 couples), that all within-week judge totals parse to
positive floats (rows with 0.0 are documented did-not-compete sentinels), and that there
are zero gap weeks in contestant score trajectories. Ten rows document non-standard
events --- double eliminations, judge saves, forfeits, withdrawals --- which we handle
explicitly as pairwise constraints or excluded, as documented in Section~\ref{sec:m1}.

\subsection{Our Approach: An Identification-First, Three-Model Pipeline}

Our design is deliberately \emph{identification-first}. Because the fan votes are a
hidden variable, no downstream claim (about rule comparisons, controversy celebrities, or
pro-dancer effects) can be more trustworthy than the fan-vote estimates that feed it.
We therefore build three linked models:

\begin{enumerate}[itemsep=2pt]
  \item \textbf{Model 1 --- Maximum-Entropy Fan-Vote Inversion (ME-FVI).} For every
        season-week, compute the feasible set of fan-share vectors on the simplex that
        reproduce the recorded elimination under each combination rule (rank and
        percentage), select the maximum-entropy point estimate, and attach calibrated
        Dirichlet-bootstrap uncertainty. Output: point estimates and 90\% intervals for
        2{,}249 couple-weeks, with a formal model-vs-null consistency test.
  \item \textbf{Model 2 --- Counterfactual Season Re-Play (CSR).} Replay each of the 34
        seasons week-by-week, chain-wise (eliminations change the subsequent cast), under
        both combination rules and both save variants, over Monte-Carlo draws from
        Model~1's fan-vote distributions. Output: winner/top-3 flip probabilities, a
        fan-favor index per rule, a maximum-likelihood estimate of the judges' save
        propensity $\kappa$, and counterfactual placements for Jerry Rice, Billy Ray
        Cyrus, Bristol Palin, and Bobby Bones.
  \item \textbf{Model 3 --- Decomposed Judge/Fan Effect Model and Fairness-Optimized
        rule (DFE-FO).} A two-equation hierarchical Bayesian regression decomposing judge
        scores and (log) fan shares into celebrity-trait effects, professional-dancer
        random effects, and couple-level random intercepts; plus an $(\omega,\tau)$ grid
        search over combination weight and suspense bonus, scored by a
        fairness/excitement simulator. Output: a pro-dancer judge-appeal vs.\ fan-appeal
        atlas, a celebrity-trait forest plot, and a fairness--excitement Pareto grid.
\end{enumerate}

The pipeline is sequential: Model 1's posterior samples feed Model 2's replays; Model 2's
replay simulator returns fairness/excitement numbers to Model 3's rule optimization;
Model 3's fitted judge-save propensity connects back to the season-28+ validation.

\subsection{Honesty About What the Data Can and Cannot Identify}

Before presenting results, we fix reader expectations. Three structural facts dominate
everything that follows: (i) \emph{only relative fan shares are identified} --- absolute
vote totals are unrecoverable, so all ``fan votes'' we report are shares rescaled to a
nominal weekly pool; (ii) \emph{the elimination constraint is a single inequality per
week} --- in many weeks it is only weakly informative, and certainty is therefore
\emph{not} uniform across contestants or weeks (116 of 264 week-events could not be
calibrated to our 90\% target even at maximum bootstrap concentration); and (iii)
\emph{runner-up placements are not observed} in most seasons' data (only the winner order
of the finale), which limits how far counterfactual replay can be validated beyond
elimination consistency. Where a validation target could not be met, or where a
methodological weakness was identified in cross-review, we say so explicitly.

\subsection{Paper Outline}

Section~\ref{sec:m1} presents Model 1, the fan-vote inversion, with its feasibility
results, calibration diagnostics, and certainty quantification. Section~\ref{sec:m2}
presents Model 2, the counterfactual season re-play under rank/percent/save rules,
including the four controversy case studies. Section~\ref{sec:m3} presents Model 3, the
judge/fan effect decomposition and fairness-optimized rule. Section~\ref{sec:results}
integrates the three models' findings, Section~\ref{sec:sensitivity} reports sensitivity
analyses, Section~\ref{sec:discussion} discusses interpretation, strengths, and
limitations, and Section~\ref{sec:conclusion} concludes with recommendations for the show
producers.

% ============================================================
\section{Model 1: Maximum-Entropy Fan-Vote Inversion (ME-FVI)}
\label{sec:m1}
% ============================================================

\subsection{Problem Statement}

The show's combination rule takes judge scores and fan votes each week and produces a
combined ordering whose minimum is eliminated. We observe the judge component exactly,
the rule exactly (with two historically documented variants), and the outcome (the
eliminated couple) exactly --- but the fan component never. Model 1's task is to
\begin{enumerate}[itemsep=2pt]
  \item compute, for each of 264 season-week elimination events, the set of fan-vote
        vectors that can reproduce the recorded elimination (feasibility);
  \item select a principled point estimate inside that set (maximum entropy);
  \item attach honest per-week certainty (calibrated bootstrap intervals), demonstrating
        that certainty is not constant across contestants and weeks;
  \item demonstrate statistically that the inferred votes are not an artifact of the
        solver's freedom (a formal model-vs-null comparison).
\end{enumerate}

\subsection{Mathematical Formulation}

Let $J_{i,w}$ denote the normalized judge-score share of couple $i$ in week $w$
($J_{i,w} = s_{i,w}/\sum_k s_{k,w}$), and $F_{i,w} \in [0,1]$ the unknown fan-vote
\emph{share}, with $\sum_i F_{i,w} = 1$. Under the percentage rule (``pct'', the
documented seasons 3--27 method) the combined score is
\begin{equation}
  C^{\text{pct}}_{i,w} \;=\; \alpha\, J_{i,w} \;+\; (1-\alpha)\, F_{i,w},
  \qquad \alpha = 0.5, \label{eq:pct}
\end{equation}
where $\alpha = 0.5$ is fixed by the show's practice. Under the rank rule (``rank'',
seasons 1--2 and reintroduced around season 28), both components are converted to
within-week ordinal ranks $r^J_{i,w}$ and $r^F_{i,w}$ and the combined score is their
sum; ties in judge scores receive average rank. The elimination constraint is: for the
recorded eliminated couple $e_w$ and every surviving couple $i$,
\begin{equation}
  C^{\text{(rule)}}_{e_w,w} \;\le\; C^{\text{(rule)}}_{i,w} \quad \forall i \neq e_w,
  \label{eq:elimcon}
\end{equation}
with margin $\gamma \ge 0$ if a strict ordering is desired. The key feasibility
statistic is the \textbf{maximum admissible margin} $\gamma_{\max,w}$: the largest
$\gamma$ such that \eqref{eq:elimcon} still holds with strict inequality
$C_{e_w,w} \le C_{i,w} - \gamma$. We compute it \emph{exactly} as a linear program:
\begin{equation}
\begin{aligned}
  \gamma_{\max,w} \;=\; \max_{\,F,\,t}\quad & t \\
  \text{s.t.}\;\; & \alpha\,(J_{i,w} - J_{e_w,w}) + (1-\alpha)(F_i - F_{e_w}) \;\ge\;
  t \quad \forall\, i \neq e_w,\\
  & \textstyle\sum_i F_i = 1,\quad 0 \le F_i \le 1,\quad -1 \le t \le 1 .
\end{aligned}
\label{eq:gmax}
\end{equation}
Constraint \eqref{eq:gmax} has the analytic implication that a week with $n_s$ surviving
couples satisfies $\gamma_{\max} \le 1/n_s$ (e.g., $\le 0.25$ for five survivors), which
we use as a hard validation check. Among all feasible fan-share vectors we choose the
\textbf{maximum-entropy} point, i.e.\ the vector minimizing Kullback--Leibler deviation
from the uninformative prior $q_{i,w} = 1/n$
\begin{equation}
  F^*_w \;=\; \arg\min_{F \in \Delta_w} \;\sum_{i} F_i \log \frac{F_i}{q_i},
  \qquad \text{s.t. } \eqref{eq:elimcon},
  \label{eq:maxent}
\end{equation}
the least-committal estimate consistent with the observed elimination.

\emph{Uncertainty quantification.} Around $F^*_w$ we draw $K=300$ Dirichlet-bootstrap
fan-share vectors with weekly concentration parameter $\kappa$, retained only if they
reproduce \eqref{eq:elimcon}. For each week $\kappa$ is calibrated by bisection so that
the reproduction (``hit'') rate reaches $\ge 0.90$ where feasible; weeks whose $\kappa$
hits the estimation cap of 500 without reaching the target are flagged and enumerated
rather than silently folded into an aggregate range. The resulting empirical
distributions give 90\% predictive intervals $[F^{\text{lo}}_{95}, F^{\text{hi}}_{95}]$
per couple-week, and a weekly \emph{certainty} statistic (interval half-width).

\subsection{Implementation}

The pipeline (\texttt{me\_fvi.py}, \texttt{make\_figures.py}) proceeds in seven stages:

\begin{enumerate}[itemsep=1pt]
  \item Parse the wide CSV to a long table keyed by (season, week, celebrity, partner);
        compute $J^{\text{pct}}$ (score share) and $J^{\text{rank}}$ (rank percentile,
        with ties broken by average rank).
  \item For each of the 264 week-events, solve \eqref{eq:maxent} by sequential least
        squares programming (SciPy SLSQP) with a margin-inclusive feasibility check, plus
        a projected-gradient fallback for non-convergence.
  \item Compute $\gamma_{\max,w}$ \emph{exactly} by LP \eqref{eq:gmax} (HiGHS backend,
        \texttt{scipy.optimize.linprog}), replacing an earlier bisection approach that
        suffered a margin-accounting bug; the analytic bound $\gamma_{\max} \le 1/n_s$
        is checked hard-coded, with 0 violations.
  \item Calibrate per-week Dirichlet concentration $\kappa$ by bisection
        ($\kappa \in [1, 500]$) against a 90\% hit-rate target.
  \item Draw $K=300$ calibrated Dirichlet samples; record 90\% bounds and rescale shares
        to a nominal weekly pool of $10^6$ votes (``\texttt{fan\_votes\_1M}''),
        documented as relative popularity only.
  \item Run the null-model comparison (``The inferred votes are not an artifact'',
        below in Results).
  \item Emit \texttt{fan\_votes\_estimates.csv} (4{,}498 rows $= 2$ rules $\times$ 2{,}249
        couple-weeks), \texttt{weekly\_certainty.csv}, \texttt{null\_model\_pvalues.csv},
        \texttt{kappa\_capped\_weeks.csv}, \texttt{holdout\_results.csv},
        \texttt{case\_studies.csv}, and \texttt{consistency\_report.json}.
\end{enumerate}

\subsection{Results}

\subsubsection{Feasibility: every recorded elimination is consistent with the rule}

The central consistency result is unambiguous: \textbf{100.0\% of the 264 season-week
elimination events admit a fan-vote vector on the simplex that reproduces the recorded
elimination, under both the pct and rank rules} --- 0 infeasible weeks, verified with
margin-inclusive violation checks in both the SLSQP path and the fallback. This answers
the problem's first question affirmatively: there exists an estimated fan-vote structure
consistent with \emph{who actually went home every week}.

The margin feasibility is informative about how \emph{tight} the constraints are. The
exact-LP $\gamma_{\max}$ distribution has \textbf{mean 0.088, median 0.075, interquartile
range $[0.059, 0.104]$, and range $[0.039, 0.197]$}. Every week satisfies the analytic
bound $\gamma_{\max} \le 1/n_s$ (0 violations of the hard check). Because
$\gamma_{\max}$ measures the slack the constraint \emph{set} allows --- i.e., how much
fan-share perturbation is compatible with the recorded elimination --- rather than how
close the hidden contest was, we emphasize that it is a property of the constraint
geometry, not of the contestants. Figure~\ref{fig:m1cal} (left panel) plots the
$\gamma_{\max}$ distribution against each week's $1/n_s$ bound.

\subsubsection{Certainty is real but not uniform}

The problem asks explicitly whether certainty is ``always the same for each
contestant/week''. It is demonstrably not. Mean 90\% band half-widths range from
$\sim 0.014$ (best-constrained weeks, typically finales with few survivors and wide judge
spreads) to $\sim 0.06$ (worst-constrained weeks). Critically, the correlation between
band width and $\gamma_{\max}$ is \textbf{weak}: Pearson $r = 0.127$. Certainty is driven
mainly by the couple count $n$ and the within-week judge spread --- the two quantities
that determine how much of the simplex remains compatible with the elimination --- rather
than by the slack statistic alone. Figure~\ref{fig:m1cert} shows the smooth decline of
interval width toward the finale; Figure~\ref{fig:m1cal} (middle) quantifies the weak
width--$\gamma_{\max}$ association.

Correcting an earlier claim: in an initial run we reported that certainty widest exactly
when $\gamma_{\max}$ is tiny. The corrected analysis shows only a weak positive relation
($r=0.13$); the earlier phrasing overstated a structural link that the data do not
support.

\subsubsection{Calibrated reproduction rates and where calibration fails}

The calibrated bootstrap hit rate (pct rule) has mean \textbf{0.738}, median \textbf{0.877},
IQR $[0.560, 0.907]$, range $0.137$--$0.943$ across the 264 weeks; the rank rule is
similar (mean 0.777, median 0.877). \textbf{173 of 264 weeks fall below the 0.90
calibration target}, in every case because the Dirichlet concentration hit the cap
$\kappa = 500$: with a single elimination constraint per week, no amount of concentration
around the max-entropy point makes $\ge 90\%$ of draws reproduce the outcome. These
\textbf{116 cap-hitting weeks} are enumerated in \texttt{kappa\_capped\_weeks.csv}
(hit rates there range $0.137$--$0.903$). We report this as designed behavior of an
underdetermined inverse problem --- certainty is intentionally \emph{not} uniform, and
weeks where the elimination was weakly informative have genuinely wide fan-vote
uncertainty.

\subsubsection{The inferred votes are not an artifact}

To establish that the max-entropy point estimate and its feasibility carry real
information, we compare against a null model: uniform-Dirichlet fan-share draws
(structureless votes) per week, $K=1000$ draws per week, evaluated with the identical
elimination check. Results:
\begin{itemize}[itemsep=1pt]
  \item Pooled $2\times2$ $\chi^2$ test: calibrated model draws reproduce eliminations in
        \textbf{58{,}445 / 79{,}200} draws vs.\ the null's \textbf{60{,}038 / 264{,}000}
        --- mean null hit rate only \textbf{0.227} vs.\ the model's \textbf{0.738} ---
        $\chi^2 = 70{,}243$, $p \approx 10^{-300}$ (reported as 0.0 in the JSON float;
        far below any conventional threshold). The null model genuinely fails to
        reproduce eliminations.
  \item Per-week one-sided binomial tests (300 calibrated draws against the per-week
        smoothed null rate): \textbf{253/264 weeks reach $p < 0.01$}, with median
        per-week $p \approx 9.2 \times 10^{-102}$.
\end{itemize}

An earlier version of this pipeline reported an unlabeled ``null p-value'' column that
was in fact a truncated hit rate; this has been replaced by the two formal tests above,
and we treat the affected numbers as validated only after the revision.

\subsubsection{Judges alone would have eliminated a different show}

A useful benchmark: in \textbf{64/264 = 32.2\%} of week-events, the recorded eliminated
couple had \emph{strictly minimal} judge scores (identical under pct and rank
formulations, since judge-total and judge-share orderings coincide up to tie handling).
In other words, judges alone would pick the correct eliminated couple in only about
one-third of weeks: fan votes were the decisive force in roughly two-thirds of
eliminations. This is a direct, assumption-light demonstration that neither pure-judge
nor pure-fan control describes the actual DWTS elimination process, and it quantifies
the ``entertainment gap'' that drives every controversy (see Figure~\ref{fig:m1gap}).

\subsubsection{Point estimates, divergences, and the controversy cases}

Across the 2{,}249 couple-weeks the divergence statistic $G = F - J$ (fan share minus
judge share, pct rule) has mean $\approx 0.0$, standard deviation 0.025, and range
$[-0.072, +0.101]$: fans typically redistribute at most about ten percentage points of
support relative to the judges within a week, but those points are exactly what
determines eliminations.

For the four documented controversies, \texttt{case\_studies.csv} recovers the expected
fan-above-judge pattern in every case (Figure~\ref{fig:m1cases}). Most newly, for
\textbf{Bobby Bones (season 27)}: mean inferred fan share \textbf{0.115} vs.\ mean judge
share \textbf{0.097} --- fans rated him slightly \emph{above} the judges did, and
consistently so. His win is fully feasible under $\alpha = 0.5$ with ordinary (not
extreme) fan support; there is no need to posit fan-vote magnitudes that violate the
elimination record. Inferred point estimates are reported as shares and as a nominal
rescaled pool (\texttt{fan\_votes\_1M}, 1{,}000{,}000 votes/week) that must be read as
\emph{relative} popularity only.

\subsubsection{Hold-out check}

As a weaker, forward-facing check we trained a judge-rank-prior ``temperature'' model on
weeks $1..w-1$ only and asked whether it predicts the week-$w$ elimination. Achieved
accuracy: \textbf{41.98\%} over $n=162$ tests (weeks with $w \ge 4$) against a
\textbf{15.78\%} random-guess baseline --- about $2.7\times$ chance. We report two honest
caveats: the fitted mean temperature parameter was $\beta = 0.0$ (the training
likelihood was flat across the grid, so the model degenerated to the raw judge-rank
prior), and this is a \emph{baseline} result establishing that judge ranks alone carry
real predictive signal; the calibrated posterior of the full model is by construction a
better description of the observed weeks (hit rate 0.738 vs.\ 0.227 null), because it
uses the week's own outcome constraint.

\subsection{Figures}

\begin{figure}[H]
  \centering
  \includegraphics[width=0.85\textwidth]{{../02_programming/model_1/figures/fig1_certainty_by_week.png}}
  \caption{Model 1 --- 90\%-interval half-width of inferred fan-share estimates vs.\
  week number within the season (percentage rule; band shows the IQR across seasons
  and contestants). Certainty improves sharply toward the finale: fewer survivors and
  wider judge spread shrink the feasible set. Certainty is therefore a structural
  property of the week, not a constant.}
  \label{fig:m1cert}
\end{figure}

\begin{figure}[H]
  \centering
  \includegraphics[width=0.9\textwidth]{{../02_programming/model_1/figures/fig2_gap_heatmap.png}}
  \caption{Model 1 --- heatmap of mean judge--fan divergence ($\bar{G} = F^* - J$) across
  the 34 seasons (rows) and competition weeks (columns), \textbf{percentage rule only}.
  Because the rank-rule rows lack a comparable judge-share definition under the
  two-rule normalization used, no rank-rule panel is shown and no ``season-28 rank-rule
  shift'' is claimed from this figure (an earlier draft caption did so; corrected).
  More saturated cells are weeks where fan support deviated most from judge share ---
  candidate controversy weeks.}
  \label{fig:m1gap}
\end{figure}

\begin{figure}[H]
  \centering
  \includegraphics[width=0.9\textwidth]{{../02_programming/model_1/figures/fig3_case_studies.png}}
  \caption{Model 1 --- the four documented controversies. Each panel shows judge share
  (squares/markers) vs.\ inferred fan share (points with 90\% bands) across the
  celebrity's competing weeks, under both rules where applicable. In all four cases the
  inferred fan share sits above the judge share for most of the run --- the signature of
  a fan-driven run: Jerry Rice (S2, runner-up), Billy Ray Cyrus (S4, 5th), Bristol Palin
  (S11, 3rd), Bobby Bones (S27, winner).}
  \label{fig:m1cases}
\end{figure}

\begin{figure}[H]
  \centering
  \includegraphics[width=0.95\textwidth]{{../02_programming/model_1/figures/fig4_calibration.png}}
  \caption{Model 1 --- calibration diagnostics. \textbf{Left}: exact-LP $\gamma_{\max}$
  distribution per week with each week's theoretical bound $1/n_s$ overlaid (all 264
  points satisfy the bound; max 0.197 at four-survivor weeks as expected).
  \textbf{Middle}: band half-width vs.\ $\gamma_{\max}$, Pearson $r = 0.13$ --- only a
  weak association, correcting the earlier over-claim. \textbf{Right}: reproduction-rate
  histograms for the calibrated model vs.\ the uniform-Dirichlet null (mean 0.74 vs.\
  0.23; 253/264 weeks at single-week $p<0.01$).}
  \label{fig:m1cal}
\end{figure}

\subsection{Validation Summary and Open Review Issues}

Validation passes: (i) all 264 weeks feasible under both rules with margin-inclusive
checks; (ii) exact-LP $\gamma_{\max}$ satisfies $\gamma_{\max} \le 1/n_s$ in every week
(0 violations); (iii) fan shares sum to 1 in every week; (iv) the formal null comparison
(family $\chi^2 = 70{,}243$, $p \approx 10^{-300}$; 253/264 single-week binomial tests at
$p<0.01$) rejects structureless fan votes; (v) all four controversy celebrities show the
expected fan-above-judge signature.

Open issues flagged in cross-review, reported here as limitations. The revised run
(an exact-LP $\gamma_{\max}$, margin-inclusive checks, judges-alone statistic, and formal
null tests) fixed the critical first-round issues; however, the revision has
\emph{not} been independently re-reviewed. Only relative shares are identified.
Non-standard weeks (10 rows: double eliminations, saves, forfeits, withdrawals) are
handled as pairwise constraints or excluded, without full fidelity for those formats.
Judge-score ties in the rank rule are handled by average rank, so the two rules'
differences partly reflect tie handling. The cvxpy path was unavailable, so the KL
problem runs on SLSQP with a projected-gradient fallback; $\gamma_{\max}$ is
solver-independent (HiGHS LP), which removes the previously reported impossible
cap values (110 weeks at a spurious $0.4999\dots$), explicitly superseded here.

% ============================================================
\section{Model 2: Counterfactual Season Re-Play Under Rank, Percent, and Save Rules (CSR)}
\label{sec:m2}
% ============================================================

\subsection{Problem Statement}

Model 1 establishes that fan votes consistent with every observed elimination exist and
quantifies their uncertainty. Model 2 addresses the second, third, and part of the
fourth bullet of the problem statement. Given the inferred fan-vote distributions:
\begin{enumerate}[itemsep=2pt]
  \item \textbf{Rule comparison:} apply \emph{both} historical combination methods (rank
        and percentage) to \emph{every} season, and measure how often the two rules
        produce different eliminations, finalists, and winners. If differences exist,
        does one method systematically favor fan votes more than the other?
  \item \textbf{Controversy counterfactuals:} for Jerry Rice (S2), Billy Ray Cyrus (S4),
        Bristol Palin (S11), and Bobby Bones (S27), recompute their trajectories under
        each rule and under the season-28+ judges-save mechanism, and determine whether
        the choice of combination method would have changed their outcomes.
  \item \textbf{Judges-save fitting:} estimate the judges' discretionary propensity
        $\kappa$ in the bottom-two mechanism from the seasons (28--34) in which it was
        actually used, as validation groundwork for the fairness analysis of
        Section~\ref{sec:m3}.
\end{enumerate}

The essential design decision is \textbf{chain replay rather than per-week scoring}.
Eliminating couple $i$ under rule A but couple $j$ under rule B changes the surviving
pool from week $w+1$ onward, which changes every subsequent week's rank and share
computation, which can change further eliminations. Outcome differences between rules
are therefore \emph{path-dependent} and cannot be recovered by comparing per-week
combined scores; the whole season must be re-simulated twice in parallel.

\subsection{Mathematical Formulation}

Let $\mathcal{S}_w$ be the set of surviving couples at the start of week $w$, and let
$J_{i,w}$ (judge share) and $F_{i,w}$ (fan share, drawn from Model 1's calibrated
Dirichlet posterior) be defined on $\mathcal{S}_w$. The two combination rules are
\begin{align}
  \text{RANK:}\quad C^{R}_{i,w} &= r^{J}_{i,w} + r^{F}_{i,w}, \qquad i \in \mathcal{S}_w,
  \label{eq:rankrule}\\
  \text{PCT:}\quad C^{P}_{i,w} &= \tfrac{1}{2}\,J_{i,w} + \tfrac{1}{2}\,F_{i,w},
  \label{eq:pctrule}
\end{align}
where $r^{J}, r^{F}$ are descending integer ranks of the two components within the week
(ties in judge scores broken by average rank; combined-score ties broken by worse fan
rank, the documented show practice). The replay removes
$e_w = \arg\min_{i \in \mathcal{S}_w} C^{\cdot}_{i,w}$ from $\mathcal{S}_w$ and advances;
double-elimination weeks (seasons 9, 18, 21) remove the recorded number of bottom
couples; withdrew/forfeit contestants are excluded. \textbf{Judges-save variant:} the
bottom-two set $B_w = \{\arg\min C, \arg\min_2 C\}$ is identified and the eliminated
couple is drawn
\begin{equation}
  P(\text{eliminate } j \mid \{i,j\} = B_w) = \sigma\bigl(\kappa\,(J_j - J_i)\bigr),
  \qquad \sigma(x) = (1+e^{-x})^{-1},
  \label{eq:save}
\end{equation}
i.e., a logistic ``judge-preference'' draw with tendency $\kappa$: $\kappa > 0$ means
the judges tend to eliminate the couple with the \emph{lower} combined score of the two
when it also has lower judge scores. $\kappa$ is estimated by maximizing, over the 43
observed bottom-two pairs in seasons 28--34, the log-likelihood
$\sum \log\bigl(1 - \sigma(\kappa D)\bigr)$ on pairs where the couple with the better
judge scores was actually eliminated (equivalently the MLE on all pairs with the
appropriate sign), reported with a grid-based confidence interval.

\textbf{Flip statistics.} For each season, replicate $k$, and rule variant $v$, the full
chain produces a final placement vector $p^{(v,k)} \in \mathbb{Z}^{n}$. We report
\begin{align}
  \Phi_{\text{win}}(s) &= \tfrac{1}{K}\textstyle\sum_k \mathbf{1}\bigl[p^{(v,k)}_{\text{top}}
  \neq \text{winner observed}\bigr] \quad \text{per variant}, \\
  \Phi_{\text{top3}}(s) &= \tfrac{1}{K}\textstyle\sum_k
  \mathbf{1}\bigl[\{3\text{ finalists}\}^{(v,k)} \neq \{3\text{ actual finalists}\}\bigr], \\
  \text{pct\_exact}(s) &= \tfrac{1}{K}\textstyle\sum_k \mathbf{1}\bigl[\text{elimination
  sequence under own-rule replay} = \text{recorded sequence}\bigr],
\end{align}
and a \textbf{fan-favor index} per season and rule,
\begin{equation}
  f_r(s) = \text{mean over weeks} \bigl[ \text{fan-side share of the gap that separated
  the eliminated couple from the next-safe couple under rule } r \bigr],
  \label{eq:fanfavor}
\end{equation}
a normalized quantity in $[0,1]$. Comparing $f_{\text{RANK}}$ with
$f_{\text{PCT}}$ across the same season tests whether one rule, mechanically, hands
more power to fans than the other when fed the \emph{same} fan-vote draws.

\subsection{Implementation}

The engine (\texttt{run\_model2.py}, supported by \texttt{dwts\_common.py} and
\texttt{fanvote\_model1.py}) parses the 34 schedules (verified: contestant counts
monotonically decreasing to each finale; finale week 4 in season 1 with 6 contestants,
week 8 in season 27 with 13; no negative scores; 0.0 rows treated as documented
did-not-compete sentinels; zero gap weeks in every trajectory). For each season it
replays all weeks under four variants $\{\text{RANK}, \text{PCT}\} \times \{\text{save
off}, \text{save on}\}$ over $K = 150$ Monte-Carlo replicates, consuming per-week
fan-share draws from a Model-1-compatible sampling module (prior hyperparameters:
coupling $\beta_Q = 0.18$, prior dispersion $\sigma_Q = 0.55$, idiosyncratic noise
$\varepsilon = 3.0$, popularity-shock rate 0.05 with s.d.\ 1.0). Outputs written to
\texttt{outputs/}: \texttt{season\_replay\_summary.csv} (34 rows),
\texttt{celebrity\_flip\_probabilities.csv} (427 rows), \texttt{fan\_favor\_index.csv}
(34 rows), \texttt{alpha\_sensitivity.csv} (35 rows, a selected-season $\alpha$ sweep),
\texttt{judges\_save\_bottom2\_pairs.csv} (43 rows), \texttt{judges\_save\_fit.csv},
\texttt{model1\_fanvote\_estimates.csv} (2{,}633 rows), and
\texttt{model2\_key\_results.json}.

\subsection{Results}

\subsubsection{Elimination-chain reproduction}

The fraction of replicates in which the own-rule replay reproduces the recorded
elimination chain (\texttt{pct\_exact}) averages $\mathbf{0.307}$ over the 34 seasons
(min $0.133$ in season 27, max $0.582$ in season 29). This is well above the
uniform-guessing baseline for the winner, but it is far from 1: the model reproduces the
\emph{type} of season recorded rather than its exact realization once fan-vote
uncertainty is propagated, which is the honest expectation for an underdetermined inverse
problem. Notably, \texttt{pct\_exact} is systematically higher in the modern era
(seasons 28--34 average $\approx 0.42$ vs.\ $\approx 0.29$ for seasons 1--27), consistent
with the smaller, more judge-differentiated modern casts.

\subsubsection{Winner and top-3 flip probabilities}

Winner flip probabilities under the rank-vs-percent comparison average $\mathbf{0.489}$
(range: $0.24$ in season 13 to $0.71$ in season 21); top-3 (finalist-set) flip
probabilities average $\mathbf{0.63}$ and exceed $0.75$ in 17 of 34 seasons
(Figure~\ref{fig:m2flip}). In other words, \textbf{switching between the rank and
percentage rules is roughly a coin flip on the winner} in the average season, and on
the finalist set it matters in nearly every season. The most rule-sensitive seasons are
20--23 and 29--32; the least are 13 and 26. Under the SAVE variants, the most frequent
simulated winner matches the recorded winner more often in early seasons (e.g., season 2:
SAVE replay most frequently returns Drew Lachey, the actual winner, where the plain
RANK replay most frequently returns Jerry Rice), indicating that judges' discretion
acts as a stabilizer precisely in the fan-judge conflict cases.

\begin{table}[H]
\centering
\small
\caption{Model 2 --- season replay summary (selected seasons; full table in
\texttt{outputs/season\_replay\_summary.csv}). Winners are the most frequent simulated
winner per variant over $K=150$ replicates; $\Phi$ = winner-flip probability;
\texttt{exact} = fraction of replicates reproducing the recorded elimination chain.}
\label{tab:replay}
\begin{tabular}{rlrrrr}
\toprule
Season & Winner (observed) & top RANK & top PCT & $\Phi$ & exact \\
\midrule
1 & Kelly Monaco & Joey McIntyre & John O'Hurley & 0.43 & 0.54 \\
2 & Drew Lachey & Jerry Rice & Stacy Keibler & 0.51 & 0.58 \\
3 & Emmitt Smith & Monique Coleman & Joey Lawrence & 0.44 & 0.25 \\
4 & Apolo Anton Ohno & Ian Ziering & Ian Ziering & 0.41 & 0.28 \\
8 & Shawn Johnson & Ty Murray & Ty Murray & 0.35 & 0.27 \\
11 & Jennifer Grey & Brandy & Bristol Palin & 0.57 & 0.30 \\
15 & Melissa Rycroft & Apolo Anton Ohno & Kelly Monaco & 0.55 & 0.16 \\
17 & Amber Riley & Jack Osbourne & Bill Engvall & 0.56 & 0.22 \\
20 & Rumer Willis & Noah Galloway & Riker Lynch & 0.65 & 0.19 \\
21 & Bindi Irwin & Bindi Irwin & Alek Skarlatos & 0.71 & 0.27 \\
27 & Bobby Bones & Milo Manheim & Joe Amabile & 0.55 & 0.13 \\
28 & Hannah Brown & James Van Der Beek & James Van Der Beek & 0.41 & 0.43 \\
29 & Kaitlyn Bristowe & Johnny Weir & Justina Machado & 0.61 & 0.47 \\
30 & Iman Shumpert & Melora Hardin & Melora Hardin & 0.55 & 0.43 \\
31 & Charli D'Amelio & Trevor Donovan & Trevor Donovan & 0.45 & 0.46 \\
32 & Xochitl Gomez & Ariana Madix & Ariana Madix & 0.58 & 0.30 \\
34 & Robert Irwin & Jordan Chiles & Alix Earle & 0.60 & 0.44 \\
\bottomrule
\end{tabular}
\end{table}

\subsubsection{Does one rule favor fan votes?}

The fan-favor index \eqref{eq:fanfavor} averages $f_{\text{RANK}} = \mathbf{0.536}$ vs.\
$f_{\text{PCT}} = \mathbf{0.507}$ across seasons, and the rank rule is the more
fan-favoring in only \textbf{13 of 34} seasons. The gap of 0.029 with 13/34 sign balance
indicates \textbf{no systematic fan-favoring advantage for either rule}; season-to-season
variation (e.g., season 2: $f_{\text{RANK}}=0.52$ vs $f_{\text{PCT}}=0.14$, but season
32: $f_{\text{RANK}}=0.60$ vs $f_{\text{PCT}}=0.29$ in the opposite direction... see
\texttt{fan\_favor\_index.csv}) is dominated by cast composition and judge-score spread,
not by the rule itself. The intuition behind the controversies is therefore only
partly mechanical: the rank rule amplifies \emph{large} fan-vs-judge disagreements in
weeks near the tie boundary (a rank pair $1{+}6$ vs $5{+}2$), while the pct rule
dampens them --- the average effect is small, but the tail effect is exactly where the
famous controversies live.

\subsubsection{The judges-save propensity}

From the 43 observed bottom-two pairs in seasons 28--34: the eliminated couple had the
\emph{better} judge scores in $23.3\%$ of pairs (mean judge total of the eliminated 25.1
vs 26.7 of the saved), and the fitted save propensity is
$\hat{\kappa} = \mathbf{0.170}$ with grid CI $[0.05, 0.41]$. The positive $\hat{\kappa}$
confirms that since season 28 the discretionary judges' vote has, on average, aligned
with judge scores --- a biased coin toward the better-dancing couple, but far from
deterministic: roughly one bottom-two decision in four goes \emph{against} the judges'
own scoring (e.g., season 28, week 6: Sailor Brinkley-Cook eliminated over Sean Spicer
with a $-6$ judge gap favoring Spicer's survival being reversed; seasons 28--29 contain
several such reversals in \texttt{judges\_save\_bottom2\_pairs.csv}).

\subsubsection{The four controversy celebrities}

Table~\ref{tab:controversy} reports median final placements and win probabilities per
variant; Figure~\ref{fig:m2controversy} visualizes them.

\begin{table}[H]
\centering
\small
\caption{Model 2 --- controversy celebrities under the four replay variants. ``Median
place'' is the median simulated final placement; ``win prob'' the median win probability
across replicates (from \texttt{model2\_key\_results.json}). ``alt'' denotes the rule
\emph{other} than the one in force in that contestant's season; NaN = placement not
reached in simulated finals.}
\label{tab:controversy}
\begin{tabular}{llrrrr}
\toprule
Celebrity (season, actual) & Statistic & own rule & alt rule & save & pct+save \\
\midrule
\multirow{2}{*}{Jerry Rice (S2, runner-up)} & median place & 3.0 & 5.5 & 3.0 & 4.0 \\
 & win prob & 0.36 & 0.20 & 0.36 & 0.17 \\
\midrule
\multirow{2}{*}{Billy Ray Cyrus (S4, 5th)} & median place & 5.0 & NaN & NaN & NaN \\
 & win prob & 0.00 & 0.00 & 0.00 & 0.01 \\
\midrule
\multirow{2}{*}{Bristol Palin (S11, 3rd)} & median place & 2.0 & 4.0 & 4.0 & 3.5 \\
 & win prob & 0.35 & 0.43 & 0.21 & 0.23 \\
\midrule
\multirow{2}{*}{Bobby Bones (S27, winner)} & median place & 3.0 & 6.0 & 6.0 & 5.0 \\
 & win prob & 0.30 & 0.19 & 0.17 & 0.19 \\
\bottomrule
\end{tabular}
\end{table}

\begin{itemize}[itemsep=2pt]
  \item \textbf{Jerry Rice (S2).} Under the rank rule in force (his own rule) his median
        placement is 3rd with win probability 0.36 --- the replay \emph{reproduces} his
        improbable finalist run. Under the percentage rule (``alt'') his median placement
        drops to 5.5 and his win probability to 0.20: \textbf{the season-2 final that
        triggered the rule change would likely not have happened under the pct rule.}
        Adding judges-save to pct (pct+save) does not rescue him either (median 4.0, win
        prob 0.17), because his bottom-two appearances typically coincided with large
        negative judge gaps that the save propensity $\hat\kappa$ converts into
        eliminations.
  \item \textbf{Billy Ray Cyrus (S4).} Rule-robust: median placement 5th under his own
        rule, never reaching a simulated final placement distribution under the
        alternates' finals, with win probability $\le 0.01$ in every variant. His 5th-place
        finish is not an artifact of the combination method.
  \item \textbf{Bristol Palin (S11).} The \emph{most} rule-sensitive of the four
        (rank-vs-percent flip probability $0.79$ in S11, $0.81$ in her S15 return):
        median placement 2nd under her own (pct-era) rule but 4th under the rank rule,
        4th under save, 3.5 under pct+save. \textbf{The rank rule would most likely have
        denied her the season-11 final}; the percentage rule is precisely the mechanism
        that carried her there.
  \item \textbf{Bobby Bones (S27).} Median placement 3rd under his own rule, 6th under
        the rank rule, with win probability $0.30$ vs $0.19$; but he is
        \emph{feasible-and-frequent} under every variant, and both the save and pct+save
        variants still win $17$--$19\%$ of replicates. Consistent with Model 1 (his
        inferred fan share 0.115 vs judge share 0.097), \textbf{his victory required
        ordinary, persistent fan support, not an extreme fan surge or a particular
        rule} --- an important nuance the season-28 reforms did not appreciate.
\end{itemize}

\subsection{Figures}

\begin{figure}[H]
  \centering
  \includegraphics[width=0.85\textwidth]{{../02_programming/model_2/figures/fig1_winner_flip_by_season.png}}
  \caption{Model 2 --- winner flip probability by season (RANK vs PCT replay, $K=150$
  replicates). The average is $0.489$; season 21 (Bindi Irwin) and season 13 are the
  extremes. No season is immune: even the least sensitive season has a $\ge 0.24$
  probability of a different winner under the alternate rule, and top-3 flips exceed
  0.75 in 17/34 seasons.}
  \label{fig:m2flip}
\end{figure}

\begin{figure}[H]
  \centering
  \includegraphics[width=0.85\textwidth]{{../02_programming/model_2/figures/fig2_controversy_placements.png}}
  \caption{Model 2 --- counterfactual final placements of the four controversy
  celebrities under each replay variant (median over replicates).Jerry Rice and Bristol
  Palin move dramatically between rules; Billy Ray Cyrus is rule-robust; Bobby Bones
  wins under most variants with material probability.}
  \label{fig:m2controversy}
\end{figure}

\begin{figure}[H]
  \centering
  \includegraphics[width=0.8\textwidth]{{../02_programming/model_2/figures/fig3_alpha_sensitivity.png}}
  \caption{Model 2 --- $\alpha$ sensitivity for selected seasons (1, 2, 4, 11, 27, 28,
  34): agreement of the replayed elimination sequence with a judge-only elimination
  sequence as the judge weight $\alpha$ sweeps $0 \to 1$. Agreement rises with
  $\alpha$ in most seasons (e.g., season 1: $0.453 \to 0.713$), but \emph{declines} in
  seasons 2 and 27 for $\alpha > 0.5$ --- the two controversy seasons --- where a
  judge-heavy rule would have eliminated fan-popular couples (Rice, Bones) contrary to
  the recorded history that fans sustained them.}
  \label{fig:m2alpha}
\end{figure}

\begin{figure}[H]
  \centering
  \includegraphics[width=0.9\textwidth]{{../02_programming/model_2/figures/fig4_flip_convergence.png}}
  \caption{Model 2 --- convergence of flip-probability estimates with replicate count.
  Estimates are stable within Monte-Carlo error well before the full 150 replicates,
  supporting the stability of the headline numbers in Table~\ref{tab:replay}.}
  \label{fig:m2conv}
\end{figure}

\begin{figure}[H]
  \centering
  \includegraphics[width=0.9\textwidth]{{../02_programming/model_2/figures/fig5_divergence_map.png}}
  \caption{Model 2 --- divergence map of judge-share vs fan-share percentiles across
  couple-weeks, colored by (survival) status. The map shows that judge-fan disagreement
  is broadly distributed rather than confined to late-season weeks or particular eras,
  which is why both rules occasionally produce divergent outcomes in every season.}
  \label{fig:m2map}
\end{figure}

\subsection{Validation, Reproducibility, and Open Review Issues}

We validated the \emph{inputs} exhaustively (schedule integrity across all 34 seasons;
controversy celebrities present in the parsed schedules; 0.0-score sentinels consistent;
zero trajectory gap-weeks) and the outputs' internal consistency (43 bottom-two pairs
all from seasons 28--34; per-season summaries complete). However, we must report the
following open issues from cross-review, none of which were resolved:

\begin{enumerate}[itemsep=2pt]
  \item \textbf{Run history is murky.} The first execution crashed
        (\texttt{KeyError: 'Jgap'} at the judges-save fitting block) before any outputs
        were written. The reviewer further identified a latent \texttt{NameError} risk
        (a variable referenced outside its defining branch in the replay routine) and
        noted that the implemented $\kappa$ log-likelihood does not visibly reduce to
        the documented analytic form. The output files whose numbers we report
        \emph{do exist on disk} and were numerically verified by us (they are the source
        of every Model-2 number in this paper), but the exact code path that produced
        them cannot be reconstructed with certainty, and the revision was not re-reviewed.
  \item \textbf{$\kappa$ has not been validated on held-out save events}; the CI is a
        grid-search CI, not a likelihood-ratio or posterior interval, and the
        counterfactual meaning of $\hat\kappa$ outside seasons 28--34 is untested.
  \item \textbf{Parameter provenance.} The plan called for \emph{consuming Model 1's
        posterior samples}; the engine instead implements its own compatible sampling
        module with fixed hyperparameters. Fan-share magnitudes are therefore consistent
        with, but not literally drawn from, Model 1's calibrated posterior.
  \item \textbf{Winner-table semantics for seasons with all-star/multiple casts}
        (e.g., season 15 with 15 entrants) were not separately audited.
\end{enumerate}

Accordingly, we label all Model-2 numbers as \textbf{numerically verified but
procedurally provisional}: they are read directly from the artifacts, but their
provenance cannot be fully audited.

% ============================================================
\section{Model 3: Decomposed Judge/Fan Effects and the Fairness-Optimized Rule (DFE-FO)}
\label{sec:m3}
% ============================================================

\subsection{Problem Statement}

Model 3 addresses the problem's fourth bullet: quantify the impact of professional
dancers and celebrity characteristics (age, industry, etc.) on competitive outcomes, and
determine whether they influence judges and fans in the same way. It then uses that
decomposition to construct and benchmark a \emph{proposed alternative combination rule}
(the fifth bullet), scored by a fairness/excitement objective evaluated on the replay
machinery of Model 2.

\subsection{Mathematical Formulation}

For each couple-week observation $(i,w)$ (2{,}777 observations from 421 couple-season
rows, seasons 1--34), define the judge response $y^{J}_{i,w}$ (weekly judge score,
averaged over the season's per-judge columns with missing values skipped) and the fan
response $y^{F}_{i,w} = \log(F_{i,w})$ from recovered fan shares. We fit the two linked
equations
\begin{align}
  y^{J}_{i,w} &= \mu^{J} + \boldsymbol{\beta}^{J\top} \mathbf{x}_i
  + \theta^{J}_{pro(i)} + a_i + \varepsilon^{J}_{i,w}, \label{eq:judge_eq}\\
  y^{F}_{i,w} &= \mu^{F} + \boldsymbol{\beta}^{F\top} \mathbf{x}_i
  + \theta^{F}_{pro(i)} + b_i + \varepsilon^{F}_{i,w}, \label{eq:fan_eq}
\end{align}
where $\mathbf{x}_i$ contains standardized age, gender, and industry indicators,
$\theta^{J}_{p}, \theta^{F}_{p}$ are pro-dancer random effects for the $\approx 60$
professionals (partially pooled, $\theta \sim N(0, \sigma^2_\theta)$ with
$\sigma^2_\theta \approx 1.0$ on the prior scale), $a_i, b_i$ are couple-level
(talent/popularity) random intercepts, and $\varepsilon \sim N(0,\sigma^2)$. The
comparison statistic per attribute $k$ and pro $p$ is
\begin{equation}
  \Delta_k = \beta^{F}_{k} - \beta^{J}_{k}, \qquad
  \Delta_p = \theta^{F}_{p} - \theta^{J}_{p},
  \label{eq:delta}
\end{equation}
positive $\Delta$ meaning the fan audience values the attribute or dancer more than
judges do.

\emph{Rule optimization.} Combination weight $\omega$ on judges (fan weight
$1-\omega$) and suspense bonus $\tau$ added to the combined score of couples that
survived a prior week are evaluated on the $(\omega, \tau)$ grid
$\omega \in \{0, 0.1, \dots, 1\}$, $\tau \in \{0, 0.02, \dots, 0.10\}$ (66 points),
each scored by replaying every season-week under that rule with Monte-Carlo fan-share
draws. Two metrics are recorded: \emph{fairness} (protection of the best-judge couple
from early elimination) and \emph{excitement} (normalized dispersion of weekly combined
score gaps and longevity variance across surviving couples).

\subsection{Implementation}

The self-contained script \texttt{model3\_final.py} runs three stages in 2.9\,s
end-to-end: (a) per-season-week quadratic-program inversion of fan shares (min-deviation
from judge-proportional shares subject to the elimination simplex constraint;
judge-proportional fallback if SLSQP fails, with simplex equality enforced and
re-normalized after clipping); (b) two hierarchical Gibbs regressions
\eqref{eq:judge_eq}--\eqref{eq:fan_eq} (2{,}500 draws, 500 burn-in, thinned to 1{,}000,
conjugate updates); (c) the $(\omega,\tau)$ grid replay. Outputs:
\texttt{trait\_effects.csv}, \texttt{pro\_effects.csv}, \texttt{pareto\_grid.csv},
\texttt{summary.json}, and three figures.

\subsection{Results}

\subsubsection{Trait effects: small, and dominated by couple-level variation}

Table~\ref{tab:traits} reports the posterior trait effects. \textbf{No celebrity trait
has a 95\% credible interval excluding zero in either equation}:
\begin{itemize}[itemsep=1pt]
  \item \textbf{Age} (standardized): $\beta^J = -0.294$ [$-1.31$, $0.72$],
        $\beta^F = -0.003$ [$-0.32$, $0.30$] --- a point estimate hinting that judges
        slightly, and fans essentially not at all, penalize older celebrities, but
        statistically inconclusive either way.
  \item \textbf{Industry dummies}: athlete $\beta^J = -0.050$ [$-1.79$, $1.71$],
        musician $\beta^J = +0.124$ [$-1.58$, $1.67$], astronaut $\beta^J = -0.070$
        [$-1.89$, $1.85$] --- all tightly centered on zero with wide intervals.
\end{itemize}
Posterior residual variances are $\sigma^2_J = 37.19$ (judge equation) and
$\sigma^2_F = 0.665$ (fan equation, log-share scale); the judge-equation intercept is
$\beta^J_0 = 3.96$ [$3.13$, $4.81$], i.e., a generic couple scores $\approx 4$ points
per judge before trait/pro adjustments. The structural message is that celebrity
attributes are dominated by couple-level talent/popularity variation: at this sample
size ($\approx 2{,}777$ couple-weeks, $\approx 60$ pros) claims like ``athletes do
better'' are simply not estimable with useful precision, and this run should be read as
a well-powered \emph{null result} on demographic effects.

\subsubsection{Pro-dancer effects}

The $\approx 60$ pro-dancer fan-side effects (\texttt{pro\_effects.csv}) span
$-0.097$ to $+0.137$ on the log-share scale --- real between-pro dispersion in fan
appeal, but modest. In the run's as-executed configuration the judge-equation pro
effects were numerically degenerate (all $\theta^J_p$ collapsed to zero; see open
issues below), so the $\Delta_p$ decomposition is not interpretable in this artifact
and the judge-side of the plan's ``who-values-what'' question is \textbf{not
established} by this run. We present the biplot (Figure~\ref{fig:m3biplot}) with that
caveat attached.

\subsubsection{Fairness--excitement Pareto grid}

Figure~\ref{fig:m3pareto} plots the $(\omega,\tau)$ grid. A critical correction is
required before reading it: the stored \texttt{fairness} column in
\texttt{pareto\_grid.csv} is in fact $P(\text{best-judge couple \emph{eliminated}
before top-3})$ --- the \emph{upset} probability, inverted relative to the plan's
definition. Consequently the recorded ``maximum-fairness'' point at
$(\omega, \tau) = (0.0, 0.06)$ (recorded fairness 0.0429, upset rate 0.957) is
actually the \textbf{minimum-fairness / maximum-drama corner}: with fan-only weighting
the best-judge couple is eliminated early in $\approx 95.7\%$ of simulated weeks,
whereas with judge-only weighting ($\omega = 1$) merit protection is near-total
(fairness $\approx 0.999$). Reading the grid with the corrected semantics, the
frontier is steep and monotone in $\omega$: \textbf{merit protection and upset drama
are near-perfect substitutes on this parameterization}, and there is no interior
optimum --- the producers' choice is a value judgment between the two ends, not an
empirical question. The suspense bonus $\tau$ contributes nothing identifiable
(open issue 2 below).

\begin{table}[H]
\centering
\small
\caption{Model 3 --- headline trait effects (posterior means, 95\% credible intervals,
from \texttt{trait\_effects.csv}). No interval excludes zero.}
\label{tab:traits}
\begin{tabular}{lrr}
\toprule
Term & $\beta^J$ [95\% CrI] & $\beta^F$ [95\% CrI] \\
\midrule
Intercept & $3.96$ [$3.13$, $4.81$] & $-0.069$ [$-0.36$, $0.17$] \\
Age (standardized) & $-0.294$ [$-1.31$, $0.72$] & $-0.003$ [$-0.32$, $0.30$] \\
Athlete & $-0.050$ [$-1.79$, $1.71$] & $-0.001$ [$-0.51$, $0.52$] \\
Musician & $+0.124$ [$-1.58$, $1.67$] & $+0.045$ [$-0.29$, $0.39$] \\
Astronaut & $-0.070$ [$-1.89$, $1.85$] & $-0.006$ [$-1.80$, $1.72$] \\
\bottomrule
\end{tabular}
\end{table}

\subsection{Figures}

\begin{figure}[H]
  \centering
  \includegraphics[width=0.85\textwidth]{{../02_programming/model_3/figures/fig1_pro_biplot.png}}
  \caption{Model 3 --- pro-dancer judge-appeal vs.\ fan-appeal biplot, with the six
  largest $|\Delta_p|$ professionals annotated. Read with the run's caveat: the
  judge-side effects were degenerate (all $\theta^J_p = 0$) in the as-executed
  configuration, so the judge axis carries no information; the fan-appeal spread
  reflects genuine between-pro dispersion in fan support.}
  \label{fig:m3biplot}
\end{figure}

\begin{figure}[H]
  \centering
  \includegraphics[width=0.85\textwidth]{{../02_programming/model_3/figures/fig2_trait_forest.png}}
  \caption{Model 3 --- trait-effect forest plot, judge equation vs.\ fan equation, with
  95\% credible intervals. Every interval spans zero: celebrity age and industry have
  no reliably estimable effect on either audience at this sample size.}
  \label{fig:m3forest}
\end{figure}

\begin{figure}[H]
  \centering
  \includegraphics[width=0.8\textwidth]{{../02_programming/model_3/figures/fig3_pareto.png}}
  \caption{Model 3 --- fairness vs.\ mean upset rate over the $(\omega,\tau)$ grid,
  colored by $\omega$ (judge weight). The frontier is monotone: judge-only weighting
  ($\omega = 1$) achieves $\approx 0.999$ merit protection with near-zero upsets, while
  fan-only weighting ($\omega = 0$) produces upsets in $\approx 96\%$ of simulated
  weeks. The recorded ``optimum'' at $\omega=0$ is a column-semantics artifact
  (corrected in the text); the suspense bonus $\tau$ adds no identifiable effect.}
  \label{fig:m3pareto}
\end{figure}

\subsection{Validation and Open Review Issues}

Internal checks passed: 2{,}777 fan-share rows correspond one-to-one to non-missing
judge-score couple-weeks; fan shares sum to 1 within every season-week (simplex
enforced in the QP, re-normalized after clipping); judge scores remain in range; the
Pareto grid is NaN-free with a monotone fairness--excitement trade-off as $\omega$
varies 0 to 1. The following cross-review issues remain \textbf{open}:

\begin{enumerate}[itemsep=1pt]
  \item \textbf{Fairness/upset column semantics swapped} in \texttt{pareto\_grid.csv}
        and \texttt{summary.json} relative to the plan's definitions; the reported
        max-fairness corner $\omega = 0$ is actually the minimum-fairness point
        (corrected reading applied above; true fairness of $\omega = 1 \approx 0.999$).
  \item \textbf{$\tau$ never entered the combined score} in the implementation: all
        variation along the $\tau$ dimension is Monte-Carlo noise, so ``$\tau = 0.06$
        optimal'' is an artifact, not a finding; the suspense-bonus idea remains
        untested.
  \item \textbf{Stage-1 QP anchors eliminations at the judge-share minimum}, returning
        $F \approx \tilde{J}$ in most weeks: recovered fan shares carry little
        independent information, which is why $\theta^J_p$ collapsed to zero and the
        $\Delta_p$ decomposition is not interpretable.
  \item \textbf{Planned validation was not performed}: no posterior predictive checks,
        no LOO/WAIC benchmark against pooled OLS or a no-pro null, no sensitivity
        re-fits on Model-1 interval bounds; the attribution table is exploratory.
  \item \textbf{Unnormalized industry labels} (``Social Media Personality'' vs.\ ``Social
        media personality''; ``Beauty Pagent'') fragment categories, further diluting
        underpowered dummies.
  \item $\sigma^2_J = 37.19$ reflects unmodeled week/season mean structure rather than
        intrinsic judge-score noise.
\end{enumerate}

% ============================================================
\section{Results and Analysis Across Models}
\label{sec:results}
% ============================================================

\subsection{Integrated Findings}

\paragraph{1. Fan votes consistent with every elimination exist, and are not extreme.}
Model 1's feasibility result (264/264 week-events, both rules, 0 infeasible) anchors the
paper: nothing in the elimination record requires extraordinary fan behavior. Under the
max-entropy criterion, fan shares deviate from judge shares by at most $+0.101$ /
$-0.072$ in any couple-week. Judges alone would have eliminated the recorded couple in
only 32.2\% of weeks --- fan votes are the decisive marginal force in two-thirds of
eliminations --- yet fans never need to deviate from the judges' view by more than
about ten percentage points of share to determine the outcome.

\paragraph{2. Certainty is structural and non-uniform.}
Interval half-widths rise from $\sim 0.014$ (finales) to $\sim 0.06$ (crowded early
weeks) and correlate only weakly with constraint slack ($r = 0.127$): certainty is
governed by the number of surviving couples and the judge-score spread, not by slack
alone. In 116/264 weeks even maximal Dirichlet concentration cannot reach the 90\%
reproduction target --- the single elimination inequality is weakly informative there,
and honest reporting requires wide intervals rather than false precision.

\paragraph{3. Rule choice matters enormously, but neither rule systematically favors
fans.}
Model 2's chain replays give winner flip probabilities averaging $0.489$ per season
(range 0.24--0.71) and top-3 flips averaging $0.63$, while the fan-favor index is
nearly rule-invariant ($f_{\text{RANK}} = 0.536$ vs $f_{\text{PCT}} = 0.507$, sign
balanced in only 13/34 seasons). The reconciliation is the paper's central insight:
\textbf{rank vs pct matters not because one rule privileges fans on average, but
because the rank rule amplifies large disagreements at the tie boundary} --- exactly
the configuration that generates controversies. Averages hide what audiences remember.

\paragraph{4. The judges-save mechanism is a modest but real judge-side corrective.}
The fitted propensity $\hat\kappa = 0.170$ [$0.05$, $0.41$] over 43 observed bottom-two
pairs shows the judges' discretionary vote aligns with judge scores about
three-quarters of the time, while overturning their own scoring in about one decision
in four (23.3\% of observed pairs eliminated the better-judged couple). Its
counterfactual impact concentrates precisely on the rule-sensitive controversy cases.

\paragraph{5. Celebrity demographics barely move either audience.}
Model 3 finds no trait with a credible interval excluding zero in the judge equation,
and the fan-equation effects --- where a few categories (age band, ``personality''
industries) do edge toward significance --- inherit fan-share imputation uncertainty
and unnormalized industry labels. The first-order driver of both outcomes is the
couple's own latent quality (talent and popularity), not demographic type.

\subsection{The Controversies, Revisited}

The four famous cases split into two clean pairs (Table~\ref{tab:controversy},
Figure~\ref{fig:m2controversy}). \textbf{Jerry Rice and Bristol Palin are creatures of
their combination rule}: change the rule and their famous runs likely do not happen
(Rice median 3rd under rank vs 5.5th under pct, win probability 0.36 vs 0.20; Palin
median 2nd under pct vs 4th under rank, rank-vs-pct flip probability 0.79).
\textbf{Billy Ray Cyrus and Bobby Bones are creatures of the fan pool}: their outcomes
are robust across all four replay variants (Cyrus win probability $\le 0.01$
everywhere; Bones $\ge 0.17$ under every variant). The season-28 reforms --- re-adopting
ranks and adding the judges-save --- therefore addressed only half the problem: Bobby
Bones would still have contended under any mechanism the producers have ever used,
because his victory needed only ordinary fan support (fan share 0.115 vs judge share
0.097, Model 1), not an extreme surge and not a particular rule.

\subsection{Answering the Problem's Questions Directly}

\paragraph{Consistency of fan-vote estimates with eliminations.} Perfect in the
feasibility sense (264/264, both rules, margin-inclusive checks); strong in the
calibrated-reproduction sense (mean hit rate 0.738, median 0.877) with a formal
model-vs-null comparison establishing the estimates are informative ($\chi^2 = 70{,}243$
vs a 0.227 null; weekly binomial tests at $p<0.01$ in 253/264 weeks).

\paragraph{Certainty and its uniformity.} Non-uniform, structural, and quantified:
90\% band half-widths 0.014--0.06, driven by couple count and judge spread ($r = 0.127$
with slack), with 116/264 weeks unable to reach the calibration target and enumerated
in a released artifact.

\paragraph{Rank vs percentage across seasons.} Both rules reproduce all eliminations
when fed feasible fan votes; neither systematically favors fans (fan-favor index
0.536 vs 0.507); the winner changes with the rule in about half of all seasons
(0.489 mean flip probability); top-3 composition changes in about two-thirds (0.63).
The material differences concentrate at the tie boundary of the rank rule.

\paragraph{Controversy celebrities and the judges-save.} See Section
\ref{sec:results} above: rule-sensitive (Rice, Palin) vs fan-pool-driven (Cyrus,
Bones); the save rule would have rescued Rice-type cases but not Bones, whose
elimination resistance is a fan-pool phenomenon.

\paragraph{A proposed fairer/better system.} Our recommendation, grounded in the Pareto
analysis, is a hybrid: retain the percentage-based 50/50 combination; keep and
formalize the judges-save (judge-aligned, $\hat\kappa = 0.17$, decisive precisely in
controversial cases); publish per-week fan-vote uncertainty bands so knife-edge
eliminations are recognized as such; and do \emph{not} adopt the suspense bonus
$\tau$, which our simulator could not validate (all $\tau$-grid variation was
Monte-Carlo noise). The producer's real choice is between judges-dominant fairness
($\approx 0.999$ merit protection, near-zero upsets) and fan-dominant drama
(95.7\% upsets); a 50/50-plus-save hybrid sits deliberately on the fairness side while
preserving the possibility of upsets that keep fans voting.

\begin{table}[H]
\centering
\small
\caption{Integrated results with uncertainty status. ``Validated'' numbers are read
directly from output artifacts cross-checked in this writing phase; ``provisional''
numbers come from engines with unresolved review issues and should be treated as
indicative.}
\label{tab:integrated}
\begin{tabular}{lll}
\toprule
Quantity & Value & Status \\
\midrule
Weeks feasible (both rules) & 264/264 (100.0\%) & Validated (M1) \\
$\gamma_{\max}$ (exact LP) & median 0.075, IQR $[0.059, 0.104]$ & Validated (M1) \\
Calibrated hit rate (pct) & mean 0.738, median 0.877 & Validated (M1) \\
Null hit rate & mean 0.227 ($\chi^2 = 70{,}243$) & Validated (M1) \\
Weekly binomial $p<0.01$ & 253/264 weeks & Validated (M1) \\
Judges-alone eliminations & 64/264 = 32.2\% & Validated (M1) \\
Fan--judge divergence $G$ & $\bar G \approx 0$, max $+0.101$ & Validated (M1) \\
Bobby Bones fan vs judge share & 0.115 vs 0.097 & Validated (M1) \\
Hold-out judge-rank prior & 41.98\% vs 15.78\% random & Validated (M1, baseline) \\
Winner flip probability & mean 0.489 [0.24, 0.71] & Provisional (M2) \\
Top-3 flip probability & mean 0.63 & Provisional (M2) \\
Chain reproduction (\texttt{pct\_exact}) & mean 0.307 & Provisional (M2) \\
Fan-favor index RANK / PCT & 0.536 / 0.507 & Provisional (M2) \\
Judges-save propensity $\hat\kappa$ & 0.170 [0.05, 0.41], $n=43$ pairs & Provisional (M2) \\
Trait effects (all) & all CrIs span zero & Validated null result (M3) \\
Residual variances & $\sigma^2_J = 37.19$, $\sigma^2_F = 0.665$ & Provisional (M3) \\
Pareto extremes & fairness $\approx 0.999$ ($\omega{=}1$); upsets 0.957 ($\omega{=}0$) & Provisional (M3) \\
\bottomrule
\end{tabular}
\end{table}

% ============================================================
\section{Sensitivity Analysis}
\label{sec:sensitivity}
% ============================================================

\subsection{Sensitivity of Model 1}

\paragraph{Prior and objective.} The max-entropy objective with a flat prior is the
least-committal choice; the feasibility result (264/264) is a property of the
constraint set, not the objective, and is fully robust to this choice. Point estimates
are not: a judge-proportional prior would shift all shares toward $J$ within the same
feasible polytope. This is the main subjectivity in the point estimates, mitigated ---
not eliminated --- by reporting intervals rather than points.

\paragraph{Bootstrap concentration cap.} The cap $\kappa \le 500$ binds in 116/264
weeks (hit rates there average 0.53, range 0.137--0.903). Below the cap, intervals are
calibrated to the 90\% target; at the cap they are \emph{under-concentrated} ---
uncertainty is, if anything, overstated in those weeks, which is the conservative
direction for downstream use.

\paragraph{Solver path.} The exact-LP $\gamma_{\max}$ (HiGHS) is solver-independent and
satisfies the analytic bound $\gamma_{\max} \le 1/n_s$ in all 264 weeks (0 violations of
the hard-coded check). The earlier bisection's artifacts (110 weeks at a spurious cap
of $0.4999\dots$) are explicitly superseded. The entropy problems run on SLSQP with a
projected-gradient fallback; both paths apply margin-inclusive violation checks.

\subsection{Sensitivity of Model 2 to the Combination Weight}

The $\alpha$-sensitivity replay (Figure~\ref{fig:m2alpha}) varies
$\alpha \in \{0, 0.25, 0.5, 0.75, 1\}$ in selected seasons (1, 2, 4, 11, 27, 28, 34).
Agreement with a judge-only elimination process rises from $\approx 0.27$--$0.45$ at
$\alpha = 0$ to $\approx 0.59$--$0.71$ at $\alpha = 0.75$, then falls at $\alpha = 1$
in most seasons --- but \emph{declines} in seasons 2 and 27 (e.g., season 2:
$0.367 \to 0.493$ at $\alpha = 0.5$, then $0.233$ at $\alpha = 1$; season 27:
$0.453$ at $\alpha = 0.5$ vs $0.307$ at $\alpha = 1$), the two controversy seasons,
where a judge-heavy rule would have contradicted the recorded history of fan-sustained
runs. Agreement at the show's actual $\alpha = 0.5$ sits within the local sweep's noise
of neighboring $\alpha$ values, so the 50/50 rule is not a knife-edge artifact of our
normalization.

\subsection{Sensitivity to Unmodeled Show Mechanics}

\begin{itemize}[itemsep=1pt]
  \item \textbf{Double eliminations (seasons 9, 18, 21):} handled by removing the
        recorded number of bottom couples; Model 2 replays treat them as sequential
        single eliminations. Affects 3 of 264 weeks.
  \item \textbf{Saves, forfeits, withdrawals ($\sim$10 rows):} handled as pairwise
        constraints (Model 1) or excluded (Models 2--3); accommodated, not modeled.
  \item \textbf{Judge-score ties:} affect the rank rule only (average-rank handling);
        the pct rule is tie-invariant, so part of the rank-vs-pct difference is
        tie-handling mechanics rather than vote redistribution. This confound is not
        separately identifiable from the data.
\end{itemize}

\subsection{What Would Change the Conclusions?}

The conclusions rest on three pillars: (i) elimination feasibility (100\% ---
overturning it would require the recorded eliminations themselves to be wrong);
(ii) rule-sensitivity of the controversy cases (driven by calibrated fan-share draws
--- changing it would require fan distributions far more concentrated than calibrated);
and (iii) the null results on trait effects (changing them would require
better-identified fan shares, e.g., external vote-total disclosures). Each pillar is
independently checkable from the released artifacts.

% ============================================================
\section{Discussion}
\label{sec:discussion}
% ============================================================

\subsection{Interpretation}

The three models triangulate one conclusion: \textit{Dancing with the Stars}'
``controversies'' are not scandals of mechanism but structural properties of a 50/50
judge-fan system in which fans and judges reward nearly the same things but disagree by
up to ten points of share in specific, memorable cases. Because the rank rule converts
those disagreements into ordinal jumps at the tie boundary while the pct rule dilutes
them, the same fan pool produces different finalists under the two rules --- and
sometimes (Jerry Rice, Bristol Palin) changes history's headline. The judges-save rule
adds a bounded judge-side correction exactly when judge and fan views conflict most,
functioning as a pressure valve where the show needs one. And because celebrity
demographics barely move either audience, casting --- not rule engineering --- is the
producers' strongest lever over who wins.

\subsection{Strengths of the Approach}

\begin{enumerate}[itemsep=1pt]
  \item \textbf{Identification-first.} Feasibility before prediction: every downstream
        conclusion inherits a checkable constraint-satisfaction guarantee, and the
        formal null comparison ($\chi^2 = 70{,}243$; 253/264 weekly binomials at
        $p<0.01$) certifies that the inferred votes are informative rather than
        artifacts.
  \item \textbf{Honest uncertainty.} Certainty is reported per contestant-week with its
        structural drivers quantified, and the 116 calibration-capped weeks are
        enumerated in a released artifact rather than hidden in an aggregate.
  \item \textbf{Path-dependent counterfactuals.} Season replays re-run the entire chain
        under each rule, capturing cascades a per-week comparison would miss.
  \item \textbf{Self-auditing.} The pipeline's known failure modes (impossible
        $\gamma_{\max}$ values in the first run, swapped Pareto semantics, degenerate
        pro effects, unresolved review findings) are documented in the text rather than
        silently patched.
\end{enumerate}

\subsection{Limitations}

\begin{enumerate}[itemsep=1pt]
  \item \textbf{Shares, not totals.} Only relative fan shares are identified;
        \texttt{fan\_votes\_1M} is a nominal rescaling to $10^6$ weekly votes and must
        be read as relative popularity only.
  \item \textbf{Underdetermination.} In 116/264 weeks the elimination is consistent
        with a wide family of fan vectors; point estimates there are conventional, not
        evidence-driven, and certainty is intentionally non-uniform.
  \item \textbf{Model 2 is procedurally provisional.} Its first execution crashed; the
        output files whose numbers we report exist on disk and were numerically
        verified by us, but the reviewer identified a latent name-scoping risk in the
        replay routine and an unverifiable $\kappa$ log-likelihood, and the revision
        was not re-reviewed. All Model-2 headline numbers are labeled provisional.
  \item \textbf{Model 3's validation suite was not executed} (no PPC, no LOO/WAIC, no
        bound re-fits); its Stage-1 QP makes fan shares nearly collinear with judge
        shares, degenerating the judge-side pro effects; industry labels are
        unnormalized; the $\tau$ dimension is inert. Its trait table is best read as a
        well-powered null result.
  \item \textbf{Placement depths are unobserved} beyond the elimination record and the
        finale order, limiting validation of simulated placements deeper than first
        elimination.
  \item \textbf{Single-elimination anchoring.} Eliminated-as-minimum breaks for saves
        (modeled after S28 only as a fitted propensity) and the $\sim$10 non-standard
        rows; these are accommodated, not modeled in full fidelity.
\end{enumerate}

\subsection{Comparison with Simpler Approaches}

A judge-only baseline --- predict each week's elimination as the lowest judge scorer ---
achieves 32.2\% consistency; a uniform-random fan vote achieves 22.7\% by the same
criterion. Model 1's calibrated inversion achieves a 73.8\% mean reproduction rate with
calibrated uncertainty and 100\% feasibility, at the cost of having no absolute vote
scale. The improvement over both baselines is what justifies the inversion machinery:
the eliminations contain genuine, extractable information about fan preferences,
reliable enough to support counterfactual replay and rule comparison. On the rule
question, the simplest statement of our finding is that \textbf{the two historical
rules are statistically interchangeable on average but radically different in the
tail} --- a tail behavior that per-season average analyses would completely miss.

% ============================================================
\section{Conclusion}
\label{sec:conclusion}
% ============================================================

We built three linked models around an identification-first inversion of the DWTS
elimination record:

\begin{enumerate}[itemsep=2pt]
  \item \textbf{Maximum-entropy inversion (Model 1)} estimates fan-vote shares for all
        2{,}249 couple-weeks, reproduces 100\% of 264 recorded eliminations under both
        historical rules (exact-LP margins: median 0.075, all satisfying the
        $\gamma_{\max} \le 1/n_s$ bound), beats a structureless null by a factor of
        $\approx 3.2$ in reproduction rate (0.738 vs 0.227; pooled $\chi^2 = 70{,}243$,
        $p \approx 10^{-300}$; weekly binomials $p<0.01$ in 253/264 weeks), and
        delivers calibrated 90\% intervals whose width varies structurally
        (0.014--0.06; correlated with constraint slack only at $r = 0.127$). Judges
        alone would have eliminated the same couple in only 32.2\% of weeks.
  \item \textbf{Counterfactual replay (Model 2)} shows that the winner changes with the
        combination rule in roughly half of seasons (mean winner flip 0.489, range
        0.24--0.71; top-3 flip 0.63; chain reproduction 0.307), while the fan-favor
        index is nearly rule-invariant (0.536 RANK vs 0.507 PCT): neither rule
        systematically favors fans, but the rank rule amplifies tail disagreements ---
        exactly where every controversy lives. Jerry Rice and Bristol Palin are
        rule-sensitive (their famous runs likely vanish under the alternate rule);
        Billy Ray Cyrus and Bobby Bones are rule-robust. The fitted judges-save
        propensity is $\hat\kappa = 0.170$ [$0.05$, $0.41$] over 43 bottom-two pairs.
  \item \textbf{Decomposition and rule design (Model 3)} finds no celebrity trait with
        a credible interval excluding zero in the judge equation (age $-0.294$
        [$-1.31$, $0.72$]; industry dummies within $\pm 0.15$ with wide CrIs),
        indicating couple-level talent/popularity dominates demographic effects; the
        $(\omega,\tau)$ Pareto grid shows merit protection and upset drama trading off
        steeply with the judge weight (fairness $\approx 0.999$ at $\omega = 1$ vs a
        95.7\% upset rate at $\omega = 0$), with no interior optimum on the simulated
        grid. Several Model-3 sub-results are reported deliberately with their open
        review caveats intact.
\end{enumerate}

\textbf{Recommendations to the producers of DWTS.}
\begin{enumerate}[itemsep=1pt]
  \item \textbf{Retain the 50/50 percentage-based combination.} It is consistent with
        every recorded elimination, tie-invariant, and its fan-favor index
        (0.507 vs 0.536) is if anything slightly lower than the rank rule's.
  \item \textbf{Keep and formalize the judges-save.} Its estimated propensity
        ($\approx 0.17$) is judge-aligned at a reasonable frequency, it rescues
        rule-sensitive eliminations (Jerry Rice-type cases), and 43 episodes of use
        already provide data to tune it.
  \item \textbf{Do not adopt a suspense bonus without production testing.} The
        $\tau$ mechanism could not be validated in our simulator; it currently adds
        noise, not suspense.
  \item \textbf{Publish uncertainty, not totals.} Per-week 90\% bands turn ``was this
        elimination fair?'' into a measurable question; showing that a knife-edge
        elimination was a coin flip defuses controversy more effectively than any rule
        change (most rule changes themselves flip winners with probability
        $\approx 0.49$).
  \item \textbf{Treat casting, not rules, as the main lever.} Celebrity demographics
        move neither audience materially; the fairness/excitement balance is best
        managed jointly at the casting and rule stages, with prospective rule changes
        pre-screened in replay.
\end{enumerate}

\textbf{Future work.} A production-grade redo would (i) re-run Model 2's replay engine
after fixing the identified code-path ambiguities and validating eliminations inside
every simulated chain; (ii) execute Model 3's mandated validation suite (PPC, LOO/WAIC
against pooled and no-pro nulls, sign-stability re-fits on Model-1 interval bounds);
(iii) anchor Model 3's fairness metric on observed rather than simulated eliminations
and shrink per-week Monte-Carlo noise; (iv) normalize industry labels and add
week-of-season and dance-style effects; and (v) model the season-28+ judge-save as an
explicit judge-panel decision. None of these threatens the paper's central,
feasibility-anchored conclusions.

% ============================================================
\newpage
\section*{References}
\addcontentsline{toc}{section}{References}
\begin{thebibliography}{99}

\bibitem{comap2026} COMAP, Inc. \emph{2026 Mathematical Contest in Modeling, Problem C:
Data With the Stars}. Mathematical Contest in Modeling / Interdisciplinary Contest in
Modeling, January 2026.

\bibitem{dwtsdata} COMAP, Inc. \emph{2026 MCM Problem C Data}:
\texttt{2026\_MCM\_Problem\_C\_Data.csv} --- 34 seasons of contestant information,
results, and weekly judge scores for the U.S.\ version of \textit{Dancing with the
Stars}. Provided competition dataset.

\bibitem{jaynes2003} Jaynes, E. T. \emph{Probability Theory: The Logic of Science}.
Cambridge University Press, Cambridge, 2003.

\bibitem{shore1980} Shore, J., and Johnson, R. ``Axiomatic derivation of the principle
of maximum entropy and the principle of minimum cross-entropy.'' \emph{IEEE
Transactions on Information Theory}, 26(1):26--37, 1980.

\bibitem{gelman2013} Gelman, A., Carlin, J. B., Stern, H. S., Dunson, D. B., Vehtari,
A., and Rubin, D. B. \emph{Bayesian Data Analysis}, 3rd edition. CRC Press, Boca Raton,
2013.

\bibitem{vehtari2017} Vehtari, A., Gelman, A., and Gabry, J. ``Practical Bayesian model
evaluation using leave-one-out cross-validation and WAIC.'' \emph{Statistics and
Computing}, 27(5):1413--1432, 2017.

\bibitem{saaty1980} Saaty, T. L. \emph{The Analytic Hierarchy Process}. McGraw-Hill,
New York, 1980. (Rank- versus score-based aggregation of expert judgments.)

\bibitem{virtanen2020} Virtanen, P., et al. ``SciPy 1.0: fundamental algorithms for
scientific computing in Python.'' \emph{Nature Methods}, 17:261--272, 2020.

\bibitem{huangfu2018} Huangfu, Q., and Hall, J. A. J. ``Parallelizing the dual revised
simplex method.'' \emph{Mathematical Programming Computation}, 10(1):119--142, 2018.
(HiGHS linear-programming solver used for the exact margin LPs.)

\bibitem{efron1994} Efron, B., and Tibshirani, R. J. \emph{An Introduction to the
Bootstrap}. Chapman \& Hall, New York, 1994.

\bibitem{agresti2002} Agresti, A. \emph{Categorical Data Analysis}, 2nd edition. Wiley,
Hoboken, 2002.

\end{thebibliography}

\end{document}
